\textbf{Phần 1. Trắc nghiệm lựa chọn}
\setcounterpageref{ex}{0}

\begin{ex} %[5D10N59]
Công thức tính vận tốc $v$ khi biết quãng đường $s$ và thời gian $t$ là
\choice
{$v = s \times t$}
{$v = t : s$}
{\True $v = s : t$}
{$v = s + t$}
\loigiai{
Vận tốc bằng quãng đường chia cho thời gian: $v = s : t$.\\
\textbf{Đáp án của câu hỏi là C}
}
\end{ex}

\begin{ex} %[5D10N59]
Đơn vị nào sau đây thường dùng để đo vận tốc của ô tô, xe máy?
\choice
{$\text{m/s}$}
{$\text{cm/phút}$}
{$\text{m/giờ}$}
{\True $\text{km/h}$}
\loigiai{
Đơn vị thông dụng để đo vận tốc ô tô, xe máy là $\text{km/h}$ (ki-lô-mét trên giờ).\\
\textbf{Đáp án của câu hỏi là D}
}
\end{ex}

\begin{ex} %[5D10H59]
Một người đi xe máy đi được quãng đường $120\text{ km}$ trong thời gian $3\text{ giờ}$. Vận tốc của người đi xe máy là
\choice
{\True $40\text{ km/h}$}
{$30\text{ km/h}$}
{$50\text{ km/h}$}
{$60\text{ km/h}$}
\loigiai{
$v = 120 : 3 = 40\text{ (km/h)}$.\\
\textbf{Đáp án của câu hỏi là A}
}
\end{ex}

\begin{ex} %[5D10H59]
Một con đà điểu chạy được quãng đường $3\,150\text{ m}$ trong $3$ phút. Vận tốc chạy của đà điểu là
\choice
{$1\,000\text{ m/phút}$}
{$1\,100\text{ m/phút}$}
{\True $1\,050\text{ m/phút}$}
{$950\text{ m/phút}$}
\loigiai{
$v = 3\,150 : 3 = 1\,050\text{ (m/phút)}$.\\
\textbf{Đáp án của câu hỏi là C}
}
\end{ex}

\begin{ex} %[5D10H59]
Một vận động viên chạy cự li $100\text{ m}$ hết $10$ giây. Vận tốc chạy của vận động viên đó là
\choice
{$1\text{ m/s}$}
{\True $10\text{ m/s}$}
{$100\text{ m/s}$}
{$0{,}1\text{ m/s}$}
\loigiai{
$v = 100 : 10 = 10\text{ (m/s)}$.\\
\textbf{Đáp án của câu hỏi là B}
}
\end{ex}

\begin{ex} %[5D10H59]
Một người đi bộ đi được quãng đường $10{,}5\text{ km}$ trong $2{,}5\text{ giờ}$. Vận tốc của người đi bộ là
\choice
{$4\text{ km/h}$}
{$4{,}5\text{ km/h}$}
{$5\text{ km/h}$}
{\True $4{,}2\text{ km/h}$}
\loigiai{
$v = 10{,}5 : 2{,}5 = 4{,}2\text{ (km/h)}$.\\
\textbf{Đáp án của câu hỏi là D}
}
\end{ex}

\begin{ex} %[5D10V59]
Một ô tô đi quãng đường $135\text{ km}$ hết $2$ giờ $15$ phút. Vận tốc của ô tô là
\choice
{$54\text{ km/h}$}
{\True $60\text{ km/h}$}
{$65\text{ km/h}$}
{$70\text{ km/h}$}
\loigiai{
Đổi $2$ giờ $15$ phút $= 2{,}25$ giờ.\\
$v = 135 : 2{,}25 = 60\text{ (km/h)}$.\\
\textbf{Đáp án của câu hỏi là B}
}
\end{ex}

\begin{ex} %[5D10V59]
Một máy bay bay được $1\,800\text{ km}$ trong $2$ giờ $30$ phút. Vận tốc của máy bay là
\choice
{$600\text{ km/h}$}
{$800\text{ km/h}$}
{$900\text{ km/h}$}
{\True $720\text{ km/h}$}
\loigiai{
Đổi $2$ giờ $30$ phút $= 2{,}5$ giờ.\\
$v = 1\,800 : 2{,}5 = 720\text{ (km/h)}$.\\
\textbf{Đáp án của câu hỏi là D}
}
\end{ex}

\begin{ex} %[5D10V59]
Một người đi xe đạp đi quãng đường $18\text{ km}$ từ $7$ giờ $15$ phút đến $8$ giờ $45$ phút. Vận tốc của người đi xe đạp là
\choice
{\True $12\text{ km/h}$}
{$10\text{ km/h}$}
{$15\text{ km/h}$}
{$18\text{ km/h}$}
\loigiai{
Thời gian đi: $8$ giờ $45$ phút $- 7$ giờ $15$ phút $= 1$ giờ $30$ phút $= 1{,}5$ giờ.\\
Vận tốc: $18 : 1{,}5 = 12\text{ (km/h)}$.\\
\textbf{Đáp án của câu hỏi là A}
}
\end{ex}

\begin{ex} %[5D10V59]
Vận tốc $36\text{ km/h}$ tương đương với bao nhiêu mét trên giây ($	ext{m/s}$)?
\choice
{$1\text{ m/s}$}
{\True $10\text{ m/s}$}
{$100\text{ m/s}$}
{$36\text{ m/s}$}
\loigiai{
$36\text{ km/h} = 36\,000\text{ m} : 3\,600\text{ s} = 10\text{ m/s}$.\\
\textbf{Đáp án của câu hỏi là B}
}
\end{ex}

\begin{ex} %[5D10V59]
Một ca nô đi xuôi dòng quãng đường $45\text{ km}$ hết $1$ giờ $30$ phút. Vận tốc của dòng nước là $3\text{ km/h}$. Vận tốc thực của ca nô khi nước yên lặng là
\choice
{$33\text{ km/h}$}
{$30\text{ km/h}$}
{\True $27\text{ km/h}$}
{$24\text{ km/h}$}
\loigiai{
Vận tốc xuôi dòng: $45 : 1{,}5 = 30\text{ (km/h)}$.\\
Vận tốc thực: $30 - 3 = 27\text{ (km/h)}$.\\
\textbf{Đáp án của câu hỏi là C}
}
\end{ex}

\begin{ex} %[5D10C59]
Một ô tô đi từ A đến B gồm một đoạn đường bằng và một đoạn đường dốc. Đoạn đường bằng dài $60\text{ km}$ ô tô đi hết $1$ giờ $12$ phút. Đoạn đường dốc dài $24\text{ km}$ ô tô đi hết $36$ phút. Vận tốc trung bình của ô tô trên cả quãng đường là
\choice
{\True $46{,}67\text{ km/h}$ (hay $\dfrac{140}{3}\text{ km/h}$)}
{$45\text{ km/h}$}
{$48\text{ km/h}$}
{$50\text{ km/h}$}
\loigiai{
Tổng quãng đường: $60 + 24 = 84\text{ (km)}$.\\
Tổng thời gian: $1$ giờ $12$ phút $+ 36$ phút $= 1$ giờ $48$ phút $= 1{,}8$ giờ.\\
Vận tốc trung bình: $84 : 1{,}8 = 46{,}67\text{ (km/h)}$.\\
\textbf{Đáp án của câu hỏi là A}
}
\end{ex}

\textbf{Phần 2. Trắc nghiệm đúng sai}
\setcounterpageref{ex}{0}

\begin{ex} %[5D10H59]
So sánh vận tốc của các phương tiện sau:
- Xe đạp: đi $15\text{ km}$ trong $1$ giờ $15$ phút.
- Xe máy: đi $42\text{ km}$ trong $1$ giờ.
- Ô tô: đi $135\text{ km}$ trong $2{,}5$ giờ.
Xét tính đúng sai:
\choiceTF
{\True Vận tốc xe đạp là $12\text{ km/h}$.}
{\True Vận tốc ô tô là $54\text{ km/h}$.}
{Vận tốc xe máy lớn hơn vận tốc ô tô.}
{\True Thứ tự vận tốc tăng dần: Xe đạp $<$ Xe máy $<$ Ô tô.}
\loigiai{
\begin{itemize}
    \item Vận tốc xe đạp: $15 : 1{,}25 = 12\text{ km/h}$. Đúng.
    \item Vận tốc ô tô: $135 : 2{,}5 = 54\text{ km/h}$. Đúng.
    \item Vận tốc xe máy ($42\text{ km/h}$) bé hơn ô tô ($54\text{ km/h}$). Khẳng định lớn hơn là Sai.
    \item $12 < 42 < 54$. Đúng.
\end{itemize}
\begin{itemize}
    \item a) \textbf{Đúng}.
    \item b) \textbf{Đúng}.
    \item c) \textbf{Sai}.
    \item d) \textbf{Đúng}.
\end{itemize}
}
\end{ex}

\begin{ex} %[5D10V59]
Một người đi xe máy từ A đến B dài $90\text{ km}$. Người đó xuất phát lúc $7$ giờ $30$ phút, dọc đường nghỉ $15$ phút và đến B lúc $10$ giờ. Xét tính đúng sai:
\choiceTF
{\True Thời gian cả đi lẫn nghỉ là $2$ giờ $30$ phút.}
{\True Thời gian thực tế đi trên đường là $2$ giờ $15$ phút $= 2{,}25$ giờ.}
{Vận tốc của người đi xe máy là $36\text{ km/h}$.}
{\True Vận tốc của người đi xe máy là $40\text{ km/h}$.}
\loigiai{
\begin{itemize}
    \item Thời gian đi + nghỉ: $10$h $- 7$h $30$p $= 2$h $30$p. Đúng.
    \item Thời gian chạy xe: $2$h $30$p $- 15$p $= 2$h $15$p $= 2{,}25$h. Đúng.
    \item Vận tốc: $90 : 2{,}25 = 40\text{ km/h} \neq 36\text{ km/h}$. Ý 3 Sai, Ý 4 Đúng.
\end{itemize}
\begin{itemize}
    \item a) \textbf{Đúng}.
    \item b) \textbf{Đúng}.
    \item c) \textbf{Sai}.
    \item d) \textbf{Đúng}.
\end{itemize}
}
\end{ex}

\textbf{Phần 3. Trả lời ngắn}
\setcounterpageref{ex}{0}

\begin{ex} %[5D10H59]
Một người chạy $800\text{ m}$ hết $4$ phút. Tính vận tốc chạy theo mét trên phút.
\shortans{200}
\loigiai{
\textbf{Đáp số:} $200$\\
$800 : 4 = 200\text{ m/phút}$.
}
\end{ex}

\begin{ex} %[5D10H59]
Một ô tô đi được $180\text{ km}$ trong $4$ giờ. Tính vận tốc của ô tô theo $\text{km/h}$.
\shortans{45}
\loigiai{
\textbf{Đáp số:} $45$\\
$180 : 4 = 45\text{ km/h}$.
}
\end{ex}

\begin{ex} %[5D10V59]
Một tàu hỏa đi quãng đường $154\text{ km}$ hết $3$ giờ $30$ phút. Tính vận tốc tàu hỏa theo $\text{km/h}$.
\shortans{44}
\loigiai{
\textbf{Đáp số:} $44$\\
$154 : 3{,}5 = 44\text{ km/h}$.
}
\end{ex}

\begin{ex} %[5D10V59]
Một con ong mật bay quãng đường $12\text{ km}$ trong $40$ phút. Tính vận tốc bay của ong mật theo $\text{km/h}$.
\shortans{18}
\loigiai{
\textbf{Đáp số:} $18$\\
$40$ phút $= \dfrac{2}{3}$ giờ $\Rightarrow 12 : \dfrac{2}{3} = 18\text{ km/h}$.
}
\end{ex}

\begin{ex} %[5D10V59]
Đổi vận tốc $72\text{ km/h}$ ra mét trên giây ($	ext{m/s}$).
\shortans{20}
\loigiai{
\textbf{Đáp số:} $20$\\
$72 : 3{,}6 = 20\text{ m/s}$.
}
\end{ex}

\begin{ex} %[5D10C59]
Một người đi xe máy từ A lúc $8$ giờ $15$ phút và đến B lúc $10$ giờ $45$ phút với vận tốc $38\text{ km/h}$. Tính quãng đường AB theo ki-lô-mét.
\shortans{95}
\loigiai{
\textbf{Đáp số:} $95$\\
Thời gian đi: $2$ giờ $30$ phút $= 2{,}5$ giờ.\\
Quãng đường: $38 \times 2{,}5 = 95\text{ km}$.
}
\end{ex}
